\chapter{2009年山东卷（理）}

\section{单选题}

\begin{problem}
集合 \(A = \{0, 2, a\}\)，\(B = \{1, a^2\}\). 若 \(A \cup B = \{0, 1, 2, 4, 16\}\)，则 \(a\) 的值为（\quad）

\choices
    {\(0\)}
    {\(1\)}
    {\(2\)}
    {\(4\)}
\end{problem}

\begin{answer}
D
\end{answer}

\begin{solution}
由 \(A\cup B=\{0,1,2,4,16\}\) 可知 \(a\in\{0,1,2,4,16\}\). 逐一代入：\(a=0\) 或 \(a=1\) 时 \(A\cup B=\{0,1,2\}\)，\(a=2\) 时 \(A\cup B=\{0,1,2,4\}\)，\(a=4\) 时 \(a^2=16\) 且 \(A\cup B=\{0,1,2,4,16\}\)，\(a=16\) 时 \(a^2=256\notin\{0,1,2,4,16\}\). 因此 \(a=4\)，故选 D.
\end{solution}

\begin{problem}
复数 \(\dfrac{3 - \i}{1 - \i}\) 等于（\quad）

\choices
    {\(1 + 2\i\)}
    {\(1 - 2\i\)}
    {\(2 + \i\)}
    {\(2 - \i\)}
\end{problem}

\begin{answer}
C
\end{answer}

\begin{solution}
分子分母同乘分母的共轭复数，得
\(\displaystyle \frac{3-\i}{1-\i}=\frac{(3-\i)(1+\i)}{(1-\i)(1+\i)}=\frac{3+3\i-\i-\i^2}{2}=\frac{4+2\i}{2}=2+\i\).
因此选 C.
\end{solution}

\begin{problem}
将函数 \( y = \sin 2x \) 的图象向左平移 \( \frac{\pi}{4} \) 个单位，再向上平移 \(1\) 个单位，所得图象的函数解析式是（\quad）

\choices
    {\(y = \cos 2x\)}
    {\(y = 2\cos^2 x\)}
    {\(y = 1 + \sin \left(2x + \frac{\pi}{4}\right)\)}
    {\(y = 2\sin^2 x\)}
\end{problem}

\begin{answer}
B
\end{answer}

\begin{solution}
图象向左平移 \(\frac{\pi}{4}\) 个单位后，函数变为 \(y=\sin\left(2x+\frac{\pi}{2}\right)=\cos2x\). 再向上平移 \(1\) 个单位，得到 \(y=1+\cos2x=2\cos^2x\)，故选 B.
\end{solution}

\begin{problem}
一空间几何体的三视图如图所示，则该几何体的体积为

\bitmapfigure[max width=0.48\linewidth]{img/2009/shandong_science/q4_fig1.png}

（\quad）

\choices
    {\(2\pi + 2\sqrt{3}\)}
    {\(4\pi + 2\sqrt{3}\)}
    {\(2\pi + \frac{2\sqrt{3}}{3}\)}
    {\(4\pi + \frac{2\sqrt{3}}{3}\)}
\end{problem}

\begin{answer}
C
\end{answer}

\begin{solution}
由三视图可知，该几何体由半径为 \(1\)、高为 \(2\) 的圆柱和一个正四棱锥组成. 俯视图中的正方形内接于直径为 \(2\) 的圆，所以正方形的对角线为 \(2\)，底面积为 \(2\). 正四棱锥的侧棱长为 \(2\)，顶点到底面中心的高为 \(\sqrt{2^2-1^2}=\sqrt3\). 因而体积为
\(\displaystyle V=\pi\cdot1^2\cdot2+\frac13\cdot2\cdot\sqrt3=2\pi+\frac{2\sqrt3}{3}\)，故选 C.
\end{solution}

\begin{problem}
已知 \(\alpha\)，\(\beta\) 表示两个不同的平面，\(m\) 为平面 \(\alpha\) 内的一条直线，则 \(\alpha \perp \beta\) 是 “\(m \perp \beta\)” 的（\quad）

\choices
    {充分不必要条件}
    {必要不充分条件}
    {充要条件}
    {既不充分也不必要条件}
\end{problem}

\begin{answer}
B
\end{answer}

\begin{solution}
若 \(m\perp\beta\)，且 \(m\subset\alpha\)，则平面 \(\alpha\) 含有一条垂直于平面 \(\beta\) 的直线，从而 \(\alpha\perp\beta\)，所以 \(\alpha\perp\beta\) 是 \(m\perp\beta\) 的必要条件. 反之，\(\alpha\perp\beta\) 时，平面 \(\alpha\) 内任意给定的直线不一定垂直于 \(\beta\). 例如取 \(\alpha:z=0\)、\(\beta:x=0\)，则 \(\alpha\perp\beta\)，但取 \(m\) 为 \(y\) 轴时，\(m\subset\alpha\) 且 \(m\subset\beta\)，所以 \(m\) 不垂直于 \(\beta\). 因此 \(\alpha\perp\beta\) 不是充分条件，故选 B.
\end{solution}

\begin{problem}
函数 \( y = \dfrac{\e^x + \e^{-x}}{\e^x - \e^{-x}} \) 的图象大致为（\quad）

\choices
    {\choicebitmap[width=0.15\paperwidth]{img/2009/shandong_science/q6_optA.png}}
    {\choicebitmap[width=0.15\paperwidth]{img/2009/shandong_science/q6_optB.png}}
    {\choicebitmap[width=0.15\paperwidth]{img/2009/shandong_science/q6_optC.png}}
    {\choicebitmap[width=0.15\paperwidth]{img/2009/shandong_science/q6_optD.png}}
\end{problem}

\begin{answer}
A
\end{answer}

\begin{solution}
函数定义域为 \(\{x\in\R\mid x\ne0\}\). 当 \(x>0\) 时，\(\e^x>\e^{-x}\)，所以 \(y>1\)；当 \(x<0\) 时，\(y<-1\). 又因为分母在 \(x=0\) 处为零，故 \(x=0\) 是竖直渐近线；当 \(x\to+\infty\) 时 \(y\to1\)，当 \(x\to-\infty\) 时 \(y\to-1\). 此外，分子为偶函数、分母为奇函数，故函数为奇函数. 图象符合这些特征的是 A.
\end{solution}

\begin{problem}
设 \(P\) 是 \(\triangle ABC\) 所在平面内的一点，\(\overrightarrow{BC} + \overrightarrow{BA} = 2\overrightarrow{BP}\)，则（\quad）

\choices
    {\(\overrightarrow{PA} +\overrightarrow{PB} = 0\)}
    {\(\overrightarrow{PC} +\overrightarrow{PA} = 0\)}
    {\(\overrightarrow{PB} +\overrightarrow{PC} = 0\)}
    {\(\overrightarrow{PA} +\overrightarrow{PB} +\overrightarrow{PC} = 0\)}
\end{problem}

\begin{answer}
B
\end{answer}

\begin{solution}
取任意原点，用位置向量表示各点，则题设为 \((\bs c-\bs b)+(\bs a-\bs b)=2(\bs p-\bs b)\)，化简得 \(2\bs p=\bs a+\bs c\). 因而 \(P\) 是线段 \(AC\) 的中点，且 \(\overrightarrow{PA}+\overrightarrow{PC}=\bs a-\bs p+\bs c-\bs p=0\)，故选 B.
\end{solution}

\begin{problem}
某工厂对一批产品进行了抽样检测. 如图是根据抽样检测后的产品净重（单位：克）数据绘制的频率分布直方图，其中产品净重的范围是\([96,106]\)，样本数据分组为\([96,98)\)，\([98,100)\)，\([100,102)\)，\([102,104)\)，\([104,106]\)，已知样本中产品净重小于\(100\text{ 克}\)的个数是\(36\)，则样本中净重大于或等于\(98\text{ 克}\)并且小于\(104\text{ 克}\)的产品的个数是（\quad）

\bitmapfigure[max width=0.40\linewidth]{img/2009/shandong_science/q8_fig1.png}

\choices
    {\(90\)}
    {\(75\)}
    {\(60\)}
    {\(45\)}
\end{problem}

\begin{answer}
A
\end{answer}

\begin{solution}
直方图中各组频率等于组距乘以相应的高. 前两组的频率为 \(2\times0.05+2\times0.10=0.30\)，所以样本总数为 \(36\div0.30=120\). 区间 \([98,104)\) 对应的三组频率为 \(2\times0.10+2\times0.15+2\times0.125=0.75\)，故所求产品数为 \(120\times0.75=90\)，选 A.
\end{solution}

\begin{problem}
设双曲线 \(\frac{x^2}{a^2} - \frac{y^2}{b^2} = 1\) 的一条渐近线与抛物线 \(y = x^2 + 1\) 只有一个公共点，则双曲线的离心率为（\quad）

\choices
    {\(\frac{5}{4}\)}
    {\(5\)}
    {\( \frac{\sqrt{5}}{2} \)}
    {\(\sqrt{5}\)}
\end{problem}

\begin{answer}
D
\end{answer}

\begin{solution}
双曲线的渐近线为 \(y=\pm\frac{b}{a}x\). 令 \(\lambda=\frac{b}{a}>0\)，将渐近线与抛物线联立，得到 \(x^2+1=\pm\lambda x\). 只有一个公共点等价于该关于 \(x\) 的二次方程判别式为零，即 \(\lambda^2-4=0\)，从而 \(\lambda=2\). 双曲线的离心率为
\(\displaystyle e=\frac{\sqrt{a^2+b^2}}{a}=\sqrt{1+\frac{b^2}{a^2}}=\sqrt5\)，故选 D.
\end{solution}

\begin{problem}
定义在 \(\R\) 上的函数 \(f(x)\) 满足 \(f(x) = \begin{cases}
\log_2(1 - x), & x \leqslant 0,\\
f(x - 1) - f(x - 2), & x > 0,
\end{cases}\) 则 \(f(2009)\) 的值为（\quad）

\choices
    {\(-1\)}
    {\(0\)}
    {\(1\)}
    {\(2\)}
\end{problem}

\begin{answer}
C
\end{answer}

\begin{solution}
记 \(u_n=f(n)\). 由已知 \(u_0=\log_2 1=0\)，\(u_{-1}=\log_2 2=1\)，所以 \(u_1=u_0-u_{-1}=-1\)，并且 \(u_n=u_{n-1}-u_{n-2}\)（\(n\geqslant1\)）. 依次计算得
\(u_1=-1\)，\(u_2=-1\)，\(u_3=0\)，\(u_4=1\)，\(u_5=1\)，\(u_6=0\). 由递推式，\(u_7=u_6-u_5=-1=u_1\)，\(u_8=u_7-u_6=-1=u_2\)，故以后每六项循环一次. 因为 \(2009=6\times334+5\)，所以 \(f(2009)=u_5=1\)，故选 C.
\end{solution}

\begin{problem}
在区间 \([-1, 1]\) 上随机取一个数 \(x\)，\(\cos \frac{\pi x}{2}\) 的值介于 \(0\) 到 \(\frac{1}{2}\) 之间的概率（\quad）

\choices
    {\(\frac{1}{3}\)}
    {\(\frac{2}{\pi}\)}
    {\(\frac{1}{2}\)}
    {\( \frac{2}{3} \)}
\end{problem}

\begin{answer}
A
\end{answer}

\begin{solution}
当 \(x\in[-1,1]\) 时，\(\frac{\pi x}{2}\in[-\frac\pi2,\frac\pi2]\). 条件 \(0\leqslant\cos\frac{\pi x}{2}\leqslant\frac12\) 等价于 \(\frac\pi3\leqslant\left|\frac{\pi x}{2}\right|\leqslant\frac\pi2\)，即 \(\frac23\leqslant|x|\leqslant1\). 满足条件的区间总长度为 \(2\times(1-\frac23)=\frac23\)，原区间长度为 \(2\)，所以概率为 \(\frac{2/3}{2}=\frac13\)，故选 A.
\end{solution}

\begin{problem}
设 \(x\)，\(y\) 满足约束条件 \(\begin{cases}
3x - y - 6 \leqslant 0,\\
x - y + 2 \geqslant 0,\\
x \geqslant 0, y \geqslant 0,
\end{cases}\) 若目标函数 \(z = ax + by\)（\(a > 0\)，\(b > 0\)）的最大值为 \(12\)，则 \(\frac{2}{a} + \frac{3}{b}\) 的最小值为（\quad）

\choices
    {\(\frac{25}{6}\)}
    {\(\frac{8}{3}\)}
    {\(\frac{11}{3}\)}
    {\(4\)}
\end{problem}

\begin{answer}
A
\end{answer}

\begin{solution}
约束条件可写为 \(y\geqslant3x-6\)、\(y\leqslant x+2\)、\(x\geqslant0\)、\(y\geqslant0\). 可行域的顶点为 \((0,0)\)、\((2,0)\)、\((4,6)\)、\((0,2)\). 在这些顶点处目标函数值分别为 \(0\)、\(2a\)、\(4a+6b\)、\(2b\). 因为 \(a\)，\(b>0\)，最大值必为 \(4a+6b\)，所以 \(4a+6b=12\)，即 \(2a+3b=6\). 由柯西不等式，
\(\displaystyle \left(\frac2a+\frac3b\right)(2a+3b)\geqslant(\sqrt2\sqrt2+\sqrt3\sqrt3)^2=25\).
因此 \(\frac2a+\frac3b\geqslant\frac{25}{6}\). 当 \(a=b=\frac65\) 时满足 \(2a+3b=6\)，且等号成立，故最小值为 \(\frac{25}{6}\)，选 A.
\end{solution}

\section{填空题}

\begin{problem}
\(\lvert 2x - 1\rvert - \lvert x - 2\rvert < 0\) 的解集为\fillinblank{}.
\end{problem}

\begin{answer}
\(\left\{x\mid -1 < x < 1\right\}\)
\end{answer}

\begin{solution}
不等式等价于 \(\lvert2x-1\rvert<\lvert x-2\rvert\). 两边均非负，可以平方，得 \((2x-1)^2<(x-2)^2\)，即 \(4x^2-4x+1<x^2-4x+4\)，所以 \(3x^2<3\)，从而 \(-1<x<1\). 解集为 \(\left\{x\mid-1<x<1\right\}\).
\end{solution}

\begin{problem}
若函数 \( f(x) = a^x - x - a\)（\(a > 0\) 且 \(a \neq 1\)）有两个零点，则实数 \(a\) 的取值范围是\fillinblank{}.
\end{problem}

\begin{answer}
\(a > 1\)
\end{answer}

\begin{solution}
当 \(0<a<1\) 时，\(f'(x)=a^x\ln a-1<0\)，函数严格递减；又 \(x\to-\infty\) 时 \(f(x)\to+\infty\)，\(x\to+\infty\) 时 \(f(x)\to-\infty\)，所以只有一个零点.

当 \(a>1\) 时，\(f''(x)=(\ln a)^2a^x>0\)，函数为严格凸函数，至多有两个零点. 并且 \(f(0)=1-a<0\)，而 \(x\to-\infty\) 和 \(x\to+\infty\) 时 \(f(x)\to+\infty\). 因而在 \((-\infty,0)\) 和 \((0,+\infty)\) 内各有一个零点，恰有两个零点. 故 \(a>1\).
\end{solution}

\begin{problem}
执行如图的程序框图，输出的 \(T=\)\fillinblank{}.

\bitmapfigure[max width=0.52\linewidth]{img/2009/shandong_science/q15_fig1.png}
\end{problem}

\begin{answer}
\(30\)
\end{answer}

\begin{solution}
初始时 \((S,T,n)=(0,0,0)\). 第一次检测 \(T>S\) 不成立，执行更新后得到 \((S,T,n)=(5,2,2)\)；第二次得到 \((10,6,4)\)；第三次得到 \((15,12,6)\)；第四次得到 \((20,20,8)\)；第五次得到 \((25,30,10)\). 此时 \(T>S\)，即 \(30>25\)，程序输出 \(T=30\).
\end{solution}

\begin{problem}
已知定义在 \(\R\) 上的奇函数 \(f(x)\)，满足 \(f(x - 4) = -f(x)\)，且在区间 \([0, 2]\) 上是增函数，若方程 \(f(x) = m\)（\(m > 0\)）在区间 \([-8, 8]\) 上有四个不同的根 \(x_1\)，\(x_2\)，\(x_3\)，\(x_4\)，则 \(x_1 + x_2 + x_3 + x_4 = \)\fillinblank{}.
\end{problem}

\begin{answer}
\(-8\)
\end{answer}

\begin{solution}
由 \(f(x-4)=-f(x)\)，将 \(x\) 换成 \(x+4\) 得 \(f(x+4)=-f(x)\)，再换一次得 \(f(x+8)=f(x)\)，所以 \(f\) 的周期为 \(8\). 又因为 \(f\) 为奇函数，结合 \(f(x+4)=-f(x)\)，有 \(f(4-x)=f(x)\). 因而若 \(u\) 是方程 \(f(x)=m\) 在 \((0,2)\) 内的根，则 \(4-u\) 也是根，且由周期性 \(u-8\) 和 \((4-u)-8\) 也是根.

由于 \(f(0)=0\)，且 \(f\) 在 \([0,2]\) 上递增，故 \(f(x)\geqslant0\)（\(0\leqslant x\leqslant2\)）. 再由 \(f(4-x)=f(x)\)，有 \(f(x)\geqslant0\)（\(2\leqslant x\leqslant4\)）. 结合 \(f(x+4)=-f(x)\) 及奇性，可知 \(f(x)\leqslant0\)（\(-4\leqslant x\leqslant0\) 或 \(4\leqslant x\leqslant8\)），所以方程 \(f(x)=m\) 的根只能位于 \([-8,-4]\cup[0,4]\). 在 \([0,2]\) 上函数递增，且 \(f(4-x)=f(x)\)，故正值水平在 \((0,4)\) 内至多有两个根；若 \(u\) 是其中一个根，则另一个为 \(4-u\)，周期性又给出 \(u-8\) 和 \((4-u)-8=-4-u\). 题设恰有四个根，所以四根可记为 \(u\)，\(4-u\)，\(u-8\)，\(-4-u\)，其和为 \(u+4-u+u-8-4-u=-8\).
\end{solution}

\section{解答题}

\begin{problem}
设函数 \( f(x) = \cos \left(2x + \frac{\pi}{3}\right) + \sin^2 x \).

\begin{enumerate}
\item 求函数 \( f(x) \) 的最大值和最小正周期；
\item 设 \(A\)，\(B\)，\(C\) 为 \(\triangle ABC\) 的三个内角，若 \(\cos B = \frac{1}{3}\)，\(f\left(\frac{C}{3}\right) = -\frac{1}{4}\)，且 \(C\) 为锐角，求 \(\sin A\).
\end{enumerate}
\end{problem}

\begin{solution}
\begin{enumerate}
\item 利用 \(\sin^2x=\frac{1-\cos2x}{2}\) 和两角和公式，得 \(\displaystyle f(x)=\frac12\cos2x-\frac{\sqrt3}{2}\sin2x+\frac{1-\cos2x}{2}=\frac12-\frac{\sqrt3}{2}\sin2x\). 因为 \(-1\leqslant\sin2x\leqslant1\)，所以 \(f(x)\) 的最大值为 \(\frac{1+\sqrt3}{2}\)，最小值为 \(\frac{1-\sqrt3}{2}\). 函数由 \(\sin2x\) 构成，且系数不为零，故其最小正周期为 \(\pi\).
\item 由第一问的化简式，\(\displaystyle -\frac14=f\left(\frac C3\right)=\frac12-\frac{\sqrt3}{2}\sin\frac{2C}{3}\)，从而 \(\sin\frac{2C}{3}=\frac{\sqrt3}{2}\). 但 \(C\) 为锐角，故 \(0<\frac{2C}{3}<\frac\pi3\)，在此区间内 \(\sin\frac{2C}{3}<\sin\frac\pi3=\frac{\sqrt3}{2}\)，矛盾. 因此原 PDF 所给条件下不存在这样的三角形，\(\sin A\) 无定义.
\end{enumerate}
\end{solution}

\begin{problem}
如图，在直四棱柱 \(ABCD - A_{1}B_{1}C_{1}D_{1}\) 中，底面 \(ABCD\) 为等腰梯形，\(AB\parallel CD\)，\(AB = 4\)，\(BC = CD = 2\)，\(AA_{1} = 2\)，\(E\)，\(E_{1}\)，\(F\) 分别是棱 \(AD\)，\(AA_{1}\)，\(AB\) 的中点.

\begin{enumerate}
\item 证明：直线 \(EE_{1}\parallel\) 平面 \(FCC_{1}\)；
\item 求二面角 \(B - FC_{1} - C\) 的余弦值.
\end{enumerate}

\bitmapfigure[max width=0.40\linewidth]{img/2009/shandong_science/q18_fig1.png}
\end{problem}

\begin{solution}
建立空间直角坐标系，取 \(D=(0,0,0)\)，\(C=(0,2,0)\)，则由等腰梯形的对称性和已知边长可取 \(A=(\sqrt3,-1,0)\)，\(B=(\sqrt3,3,0)\). 于是 \(C_1=(0,2,2)\)，\(E=(\frac{\sqrt3}{2},-\frac12,0)\)，\(E_1=(\sqrt3,-1,1)\)，\(F=(\sqrt3,1,0)\).

\begin{enumerate}
\item 设 \(F_1\) 为 \(A_1B_1\) 的中点，则 \(F_1=(\sqrt3,1,2)\). 有 \(\overrightarrow{CF_1}=(\sqrt3,-1,2)\)，而 \(\overrightarrow{EE_1}=(\frac{\sqrt3}{2},-\frac12,1)=\frac12\overrightarrow{CF_1}\). 又 \(C\)，\(F\)，\(F_1\)，\(C_1\) 构成平行四边形，故 \(CF_1\subset\text{平面 }FCC_1\). 平面 \(FCC_1\) 的方程为 \(x+\sqrt3y-2\sqrt3=0\)，而点 \(E\) 代入左端所得为 \(-2\sqrt3\ne0\)，所以点 \(E\) 不在平面 \(FCC_1\) 内. 因而 \(EE_1\parallel CF_1\)，且 \(EE_1\) 不与平面 \(FCC_1\) 相交，从而 \(EE_1\parallel\text{平面 }FCC_1\).
\item 平面 \(BFC_1\) 的两个方向向量可取 \(\overrightarrow{FB}=(0,2,0)\) 和 \(\overrightarrow{FC_1}=(-\sqrt3,1,2)\)，所以其法向量可取 \(\bs n_1=(2,0,\sqrt3)\). 平面 \(CFC_1\) 的两个方向向量可取 \(\overrightarrow{FC}=(-\sqrt3,1,0)\) 和 \(\overrightarrow{FC_1}=(-\sqrt3,1,2)\)，所以其法向量可取 \(\bs n_2=(1,\sqrt3,0)\). 二面角取锐角时，
\(\displaystyle \cos\theta=\frac{|\bs n_1\cdot\bs n_2|}{|\bs n_1||\bs n_2|}=\frac{2}{\sqrt7\cdot2}=\frac{\sqrt7}{7}\).
\end{enumerate}
\end{solution}

\begin{problem}
在某校组织的一次篮球定点投篮训练中，规定每人最多投\(3\text{ 次}\)；在 \(A\) 处每投进一球得\(3\text{ 分}\)，在 \(B\) 处每投进一球得\(2\text{ 分}\)；如果前两次得分之和超过\(3\text{ 分}\)即停止投篮，否则投第\(3\text{ 次}\). 某同学在 \(A\) 处的命中率 \(q_1\) 为 \(0.25\)，在 \(B\) 处的命中率为 \(q_2\)，该同学选择先在 \(A\) 处投一球，以后都在 \(B\) 处投，用 \(\xi\) 表示该同学投篮训练结束后所得的总分，其分布列为

\begin{center}
\begin{tabular}{c|ccccc}
\(\xi\) & \(0\) & \(2\) & \(3\) & \(4\) & \(5\) \\
\hline
\(p\) & \(0.03\) & \(p_1\) & \(p_2\) & \(p_3\) & \(p_4\)
\end{tabular}
\end{center}

\begin{enumerate}
\item 求 \(q_2\) 的值；
\item 求随机变量 \(\xi\) 的数学期望 \(E(\xi)\)；
\item 试比较该同学选择都在 \(B\) 处投篮得分超过\(3\)分与选择上述方式投篮得分超过\(3\)分的概率的大小.
\end{enumerate}
\end{problem}

\begin{solution}
设在 \(A\) 处投中为事件 \(A\)，在 \(B\) 处投中为事件 \(B\)，各次投篮相互独立.

\begin{enumerate}
\item 得分为 \(0\) 的唯一情形是三次均未投中，所以 \(0.03=(1-q_1)(1-q_2)^2=0.75(1-q_2)^2\). 因而 \((1-q_2)^2=0.04\). 因为 \(0\leqslant q_2\leqslant1\)，所以 \(q_2=0.8\).
\item 代入 \(q_1=0.25\)、\(q_2=0.8\)，各分数的概率为
\(P(\xi=2)=0.75\times0.2\times0.8+0.75\times0.8\times0.2=0.24\)，
\(P(\xi=3)=0.25\times0.2\times0.2=0.01\)，
\(P(\xi=4)=0.75\times0.8\times0.8=0.48\)，
\(P(\xi=5)=0.25\times0.8+0.25\times0.2\times0.8=0.24\). 因而
\(\displaystyle E(\xi)=0\times0.03+2\times0.24+3\times0.01+4\times0.48+5\times0.24=3.63\).
\item 若三次都在 \(B\) 处投篮，则得分超过 \(3\) 分等价于至少投中两球，概率为 \(\displaystyle \mathrm C_3^2(0.8)^2(0.2)+(0.8)^3=0.896\). 按题设方式投篮时，得分超过 \(3\) 分的概率为 \(P(\xi=4)+P(\xi=5)=0.48+0.24=0.72\). 因为 \(0.896>0.72\)，所以都在 \(B\) 处投篮的概率较大.
\end{enumerate}
\end{solution}

\begin{problem}
等比数列 \(\{a_{n}\}\) 的前 \(n\) 项和为 \(S_{n}\). 已知对任意的 \(n \in \N^*\)，点 \((n, S_{n})\) 均在函数 \(y = b^{x} + r\)（\(b > 0\) 且 \(b \neq 1\)，\(b\)，\(r\) 均为常数）的图象上.

\begin{enumerate}
\item 求 \(r\) 的值；
\item 当 \(b = 2\) 时，记 \(b_{n} = 2(\log_{2} a_{n} + 1)\)（\(n \in \N^*\)）. 证明：对任意的 \(n \in \N^*\)，不等式 \(\frac{b_{1} + 1}{b_{1}} \cdot \frac{b_{2} + 1}{b_{2}} \cdots \frac{b_{n} + 1}{b_{n}} > \sqrt{n + 1}\) 成立.
\end{enumerate}
\end{problem}

\begin{solution}
\begin{enumerate}
\item 由题设 \(S_n=b^n+r\). 因而 \(a_1=S_1=b+r\)，\(a_2=S_2-S_1=b^2-b=b(b-1)\)，且 \(a_3=S_3-S_2=b^3-b^2=b^2(b-1)=ba_2\). 因为 \(\{a_n\}\) 是等比数列且 \(a_2\ne0\)，其公比为 \(a_3/a_2=b\)，所以 \(a_2=ba_1\). 于是 \(b(b-1)=b(b+r)\). 由于 \(b>0\)，可约去 \(b\)，得到 \(r=-1\).
\item 由第一问，\(S_n=2^n-1\)，所以 \(a_n=S_n-S_{n-1}=2^{n-1}\). 从而 \(b_n=2(\log_2 2^{n-1}+1)=2n\). 记左端乘积为 \(P_n\)，则
\(\displaystyle P_n^2=\prod_{k=1}^n\left(\frac{2k+1}{2k}\right)^2\).
对每个正整数 \(k\)，有
\(\displaystyle \left(\frac{2k+1}{2k}\right)^2-\frac{k+1}{k}=\frac{1}{4k^2}>0\)，所以
\(\displaystyle P_n^2>\prod_{k=1}^n\frac{k+1}{k}=n+1\). 因为 \(P_n>0\)，故 \(P_n>\sqrt{n+1}\)，不等式得证.
\end{enumerate}
\end{solution}

\begin{problem}
两县城 \(A\) 和 \(B\) 相距 \(20 \mathrm{~km}\)，现计划在两县城外以 \(AB\) 为直径的半圆弧 \(\widehat{AB}\) 上选择一点 \(C\) 建造垃圾处理厂，其对城市的影响度与所选地点到城市的距离有关，对城 \(A\) 和城 \(B\) 的总影响度为城 \(A\) 与城 \(B\) 的影响度之和. 记 \(C\) 点到城 \(A\) 的距离为 \(x \mathrm{~km}\)，建在 \(C\) 处的垃圾处理厂对城 \(A\) 和城 \(B\) 的总影响度为 \(y\). 统计调查表明：垃圾处理厂对城 \(A\) 的影响度与所选地点到城 \(A\) 的距离的平方成反比. 比例系数为 \(4\). 对城 \(B\) 的影响度与所选地点到城 \(B\) 的距离的平方成反比，比例系数为 \(k\). 当垃圾处理厂建在 \(\widehat{AB}\) 的中点时，对城 \(A\) 和城 \(B\) 的总影响度为 \(0.065\).

\begin{enumerate}
\item 将 \(y\) 表示成 \(x\) 的函数；
\item 讨论（1）中函数的单调性，并判断弧 \(\widehat{AB}\) 上是否存在一点，使建在此处的垃圾处理厂对城 \(A\) 和城 \(B\) 的总影响度最小？若存在，求出该点到城 \(A\) 的距离；若不存在，说明理由.
\end{enumerate}
\end{problem}

\begin{solution}
\begin{enumerate}
\item 因为 \(AB\) 是直径，故 \(\angle ACB=90^\circ\)，从而 \(BC^2=AB^2-AC^2=400-x^2\)，且 \(0<x<20\). 在半圆弧 \(\widehat{AB}\) 的中点处，\(AC=BC\)，结合 \(AC^2+BC^2=AB^2=400\)，得 \(AC=BC=10\sqrt2\). 所以 \(\frac4{(10\sqrt2)^2}+\frac{k}{(10\sqrt2)^2}=0.065\)，解得 \(k=9\). 因此
\(\displaystyle y=\frac4{x^2}+\frac{9}{400-x^2}\)，其中 \(0<x<20\).
\item 求导得
\(\displaystyle y'=-\frac8{x^3}+\frac{18x}{(400-x^2)^2}=\frac{18x^4-8(400-x^2)^2}{x^3(400-x^2)^2}\).
因为 \(x>0\)，令 \(y'=0\)，得 \(18x^4=8(400-x^2)^2\)，即 \(3x^2=2(400-x^2)\)，所以 \(x^2=160\)，即 \(x=4\sqrt{10}\). 当 \(0<x<4\sqrt{10}\) 时，\(18x^4<8(400-x^2)^2\)，故 \(y'<0\)；当 \(4\sqrt{10}<x<20\) 时，\(18x^4>8(400-x^2)^2\)，故 \(y'>0\). 因此函数在 \(0<x<4\sqrt{10}\) 上单调递减，在 \(4\sqrt{10}<x<20\) 上单调递增，故弧上存在唯一最小值点，其到城 \(A\) 的距离为 \(4\sqrt{10}\mathrm{~km}\).
\end{enumerate}
\end{solution}

\begin{problem}
设椭圆 \(E: \frac{x^2}{a^2} + \frac{y^2}{b^2} = 1\)（\(a > 0\)，\(b > 0\)）过 \(M(2, \sqrt{2})\)，\(N(\sqrt{6}, 1)\) 两点，\(O\) 为坐标原点.

\begin{enumerate}
\item 求椭圆 \(E\) 的方程；
\item 是否存在圆心在原点的圆，使得该圆的任意一条切线与椭圆 \(E\) 恒有两个交点 \(A\)，\(B\)，且 \(\overrightarrow{OA} \perp \overrightarrow{OB}\)？若存在，写出该圆的方程，并求 \(|AB|\) 的取值范围；若不存在，说明理由.
\end{enumerate}
\end{problem}

\begin{solution}
\begin{enumerate}
\item 将两点代入椭圆方程，得 \(\frac4{a^2}+\frac2{b^2}=1\) 和 \(\frac6{a^2}+\frac1{b^2}=1\). 令 \(u=\frac1{a^2}\)，\(v=\frac1{b^2}\)，则 \(4u+2v=1\)，\(6u+v=1\). 解得 \(u=\frac18\)，\(v=\frac14\)，所以椭圆方程为 \(\displaystyle \frac{x^2}{8}+\frac{y^2}{4}=1\).
\item 设所求圆的半径为 \(r\)，先取非竖直切线 \(l:y=kx+m\). 将其代入椭圆方程，交点横坐标 \(x_1\)，\(x_2\) 满足 \(\displaystyle (1+2k^2)x^2+4kmx+2m^2-8=0\). 由韦达定理，结合 \(y_i=kx_i+m\)，有
\(\displaystyle x_1x_2=\frac{2m^2-8}{1+2k^2}\)，
\(\displaystyle y_1y_2=\frac{m^2-8k^2}{1+2k^2}\)，
从而
\(\displaystyle \overrightarrow{OA}\cdot\overrightarrow{OB}=x_1x_2+y_1y_2=\frac{3m^2-8k^2-8}{1+2k^2}\).
因直线与圆相切，\(m^2=r^2(1+k^2)\). 要使上式对任意 \(k\) 恒为零，代入并整理得 \(3r^2=8\)，故 \(r^2=\frac83\). 对上述非竖直切线，代入 \(m^2=\frac83(1+k^2)\) 后，交点方程的判别式为 \(\displaystyle \Delta=\frac{32(1+4k^2)}{3}>0\)，故均为两个不同交点. 对于竖直切线 \(x=\pm r\)，其与椭圆的两个交点的纵坐标为 \(y=\pm2\sqrt{1-\frac{r^2}{8}}=\pm2\sqrt{\frac23}\)，所以 \(\overrightarrow{OA}\cdot\overrightarrow{OB}=r^2-\frac83=0\)，同样满足正交条件. 因而确实存在所求圆，方程为 \(x^2+y^2=\frac83\).

\(\displaystyle |AB|^2=(1+k^2)(x_1-x_2)^2=\frac{32(1+k^2)(1+4k^2)}{3(1+2k^2)^2}\). 令 \(t=k^2\geqslant0\)，则需研究 \(\displaystyle h(t)=\frac{(1+t)(1+4t)}{(1+2t)^2}\). 又有 \(\displaystyle h(t)-1=\frac{t}{(1+2t)^2}\geqslant0\)，且
\(\displaystyle h'(t)=\frac{1-2t}{(1+2t)^3}\)，所以 \(h(t)\) 在 \(t=\frac12\) 处取得最大值 \(\frac98\)，最小值为 \(1\). 竖直切线给出同样的最小值，故
\(\displaystyle \frac{32}{3}\leqslant|AB|^2\leqslant12\)，即 \(\displaystyle \frac{4\sqrt6}{3}\leqslant|AB|\leqslant2\sqrt3\).
\end{enumerate}
\end{solution}
